\documentclass[11pt]{article}
\usepackage[T1]{fontenc}
\usepackage[utf8]{inputenc}
\usepackage{fourier}
\usepackage[scaled=0.875]{helvet}
\renewcommand{\ttdefault}{lmtt}
\usepackage{amsmath,amssymb}
\usepackage{makeidx}
\makeindex
\usepackage{fancybox}
\usepackage[normalem]{ulem}
\usepackage{enumitem}
\usepackage{pifont}
\newcommand{\textding}[1]{\text{\ding{#1}}}
\usepackage{lscape}
%\usepackage{lastpage}
\usepackage{tabularx,booktabs,bookmark}
\usepackage{scratch3}
\usepackage{textcomp}
\usepackage{enumitem}
\usepackage{multicol}
\usepackage{multirow}
%\usepackage{mdframed}
%\definecolor{scrmovedddd}{HTML}{3373cc}%
\newcommand{\euro}{\eurologo{}}
%Tapuscrit : Denis Vergès
\usepackage{pst-eucl}
%\usepackage{diagbox}% à mettre après pst-eucl
\usepackage{graphicx,pstricks,pst-plot,pst-grad,pst-node,pst-text,pstricks-add}
\usepackage{pstricks-add}
\usepackage{tikz}
\usetikzlibrary{arrows.meta, calc, angles, quotes}
\usepackage{pgfplots}
\pgfplotsset{compat=1.18}
\definecolor{gris}{gray}{0.85}
\usetikzlibrary{patterns,calc,decorations.pathmorphing}
\newcommand{\titreexo}[1]{\par\noindent\textbf{#1}\par\vspace{0.2em}}
\newcommand{\R}{\mathbb{R}}
\newcommand{\N}{\mathbb{N}}
\newcommand{\D}{\mathbb{D}}
\newcommand{\Z}{\mathbb{Z}}
\newcommand{\Q}{\mathbb{Q}}
\newcommand{\C}{\mathbb{C}}
\usepackage[left=3.5cm, right=3.5cm, top=2.2cm, bottom=2cm]{geometry}
\setlength\headheight{13.7pt}
\newcommand{\vect}[1]{\overrightarrow{\,\mathstrut#1\,}}
\renewcommand{\theenumi}{\textbf{\arabic{enumi}}}
\renewcommand{\labelenumi}{\textbf{\theenumi.}}
\renewcommand{\theenumii}{\textbf{\alph{enumii}}}
\renewcommand{\labelenumii}{\textbf{\theenumii.}}
\def\Oij{$\left(\text{O}~;~\vect{\imath},~\vect{\jmath}\right)$}
\def\Oijk{$\left(\text{O}~;~\vect{\imath},~\vect{\jmath},~\vect{k}\right)$}
\def\Ouv{$\left(\text{O}~;~\vect{u},~\vect{v}\right)$}
% Tableau QCM
\newenvironment{qcm}{%
\begin{center}\renewcommand\arraystretch{1.4}
\begin{tabular}{|>{\centering\arraybackslash}p{3.2cm}|>{\centering\arraybackslash}p{3.2cm}|>{\centering\arraybackslash}p{3.2cm}|>{\centering\arraybackslash}p{3.2cm}|}
\hline
}{%
\\ \hline
\end{tabular}
\renewcommand\arraystretch{1}
\end{center}
}
%\newcommand{\arc}[1]{\widehat{#1}}
\usepackage{fancyhdr}
\usepackage[french]{babel}
\usepackage[np]{numprint}
\usepackage{colortbl}
%\usepackage[colorlinks=true,linkcolor=blue,citecolor=blue,urlcolor=blue]{hyperref}
%\hypersetup{%
%pdfauthor = {APMEP},
%pdfsubject = {Brevet},
%pdftitle = {Année 2026},
%allbordercolors = white,
%pdfstartview=FitH}
\begin{document}
\setlength\parindent{0mm}
%\marginpar{\rotatebox{90}{\textbf{A. P{}. M. E. P{}.}}}
\lhead{\small L'année 2026}
\rhead{\small \textbf{A. P{}. M. E. P{}.}}
\rfoot{\small Asie}
\lfoot{\small 16 juin 2026}
\pagestyle{fancy}
\thispagestyle{empty}
\begin{center} {\Large \textbf{\decofourleft~Diplôme national du brevet Asie ~\decofourright\\[7pt] 16 juin 2026}}
\end{center}

%\section*{Exercice 1 :\hfill 20 points}

%\begin{center}
%  {\large\textbf{DIPLÔME NATIONAL DU BREVET}}\\[4pt]
%  {\large\textbf{SESSION 2026}}\\[20pt]
%  \fbox{\parbox{0.75\textwidth}{\centering
%    \vspace{12pt}
%    {\LARGE\textbf{MATHÉMATIQUES}}\\[10pt]
%    {\large\textbf{Série générale}}\\[12pt]
%    \textbf{Durée : 2 h 00} \hfill \textbf{Coefficient : 2}\\[8pt]
%  }}
%\end{center}
%
%\vspace{12pt}
%\begin{center}
%  Dès que le sujet vous est remis, assurez-vous qu'il est complet.\\
%  Ce sujet comporte 8 pages numérotées de la page 1/8 à la page 8/8.
%\end{center}

%\vspace{12pt}
\begin{center}
\begin{tabular}{|p{0.60\textwidth}|p{0.25\textwidth}|}\hline
\textbf{Partie 1 – Automatismes} \newline 20 min (calculatrice interdite) & \textbf{6 points} \\ \hline
\textbf{Partie 2 – Raisonnement et résolution de problèmes} \newline 1 h 40 (calculatrice autorisée) & \textbf{14 points} \\ \hline
\end{tabular}
\end{center}

\vspace{12pt}

\begin{center}
\textbf{À l'issue de la partie 1, les copies sont ramassées.}
\end{center}

\medskip

L'usage de la calculatrice avec mode examen actif ou sans mémoire \og type collège \fg est \textbf{interdit} pour la partie 1 et \textbf{autorisé} pour la partie 2.

L'utilisation du dictionnaire est interdite.

\bigskip

\section*{Partie 1 – Automatismes – 6 points – 20 minutes}

Pour chaque question, recopier sur la copie son numéro et la réponse correspondante.

Pour cette partie, aucune justification n'est demandée.

Pour les questions à choix multiple, une seule réponse est exacte.

\titreexo{Question 1}
L'écriture scientifique du nombre \np{45310} est :

\begin{qcm}
\textbf{Réponse A} & \textbf{Réponse B} & \textbf{Réponse C} & \textbf{Réponse D} \\
\hline
$45,31 \times 10^{3}$ & $4,531 \times 10^{4}$ & $4,531 \times 10^{-4}$ & $4531 \times 10^{1}$
\end{qcm}

\titreexo{Question 2}
Une forme développée de l'expression $(4x-3)(4x+3)$ est :

\begin{qcm}
\textbf{Réponse A} & \textbf{Réponse B} & \textbf{Réponse C} & \textbf{Réponse D} \\
\hline
$4x^2 - 9$ & $16x^2+9$ & $16x^2-9$ & $8x^2-6$
\end{qcm}

\titreexo{Question 3}
Un pavé droit a pour dimensions : 4,5 cm de long, 4 cm de large, 10 cm de haut. Le volume de ce pavé est de :

\begin{qcm}
\textbf{Réponse A} & \textbf{Réponse B} & \textbf{Réponse C} & \textbf{Réponse D} \\
\hline
$180~\text{cm}^3$ & $170~\text{cm}^3$ & $160,5~\text{cm}^3$ & $18,5~\text{cm}^3$
\end{qcm}

\titreexo{Question 4}
On considère les nombres suivants et on s'intéresse à leur divisibilité par 9.

\[N = \np{2025} \quad \text{et} \quad P = \np{2026}.\]

\begin{qcm}
\textbf{Réponse A} & \textbf{Réponse B} & \textbf{Réponse C} & \textbf{Réponse D} \\
\hline
\small $N$ et $P$ sont tous les deux divisibles par 9 & $N$ est divisible par 9 mais $P$ ne l'est pas & $P$ est divisible par 9 mais $N$ ne l'est pas & Aucun des deux n'est divisible par 9
\end{qcm}

\titreexo{Question 5}
Une personne a couru 9~km en 45~minutes.

Quelle est sa vitesse moyenne en km/h ?

\titreexo{Question 6}

\begin{minipage}{0.68\textwidth}
Une roue de la fortune est utilisée pour faire gagner des cadeaux. La roue est divisée en 10 secteurs de tailles égales, avec les gains suivants : des stylos, des porte-clés, des casques audios ou un smartphone.

Un joueur tourne la roue une seule fois.

Quelle est la probabilité que le joueur gagne un casque audio ?
\end{minipage}%
\hfill
\begin{minipage}{0.28\textwidth}
\centering
\begin{tikzpicture}[scale=1.1]
\foreach \i in {0,...,9} {
  \draw (0,0) -- (90+36*\i:1.6);
}
\draw (0,0) circle (1.6);
\node at (108:1.1) {\tiny stylo};
\node at (72:1.1) {\tiny casque};
\node at (36:1.1) {\tiny porte-clé};
\node at (0:1.1) {\tiny smartphone};
\node at (-36:1.1) {\tiny porte-clé};
\node at (-72:1.1) {\tiny stylo};
\node at (-108:1.1) {\tiny stylo};
\node at (-144:1.1) {\tiny casque};
\node at (180:1.1) {\tiny porte-clé};
\node at (144:1.1) {\tiny stylo};
\end{tikzpicture}
\end{minipage}

\titreexo{Question 7}
Un article coûte 60 \euro. Calculer son nouveau prix après une baisse de 10\,r

\%.

\titreexo{Question 8}

\begin{minipage}{0.62\textwidth}
Quelle est la mesure de l'angle $\widehat{\text{BAC}}$ ?
\end{minipage}%
\hfill
\begin{minipage}{0.34\textwidth}
\centering
\begin{tikzpicture}[scale=0.8]
\coordinate (C) at (0,0);
\coordinate (A) at (3,0);
\coordinate (B) at (0,3.5);
\draw (B) -- (C) -- (A) -- (B);
\draw (0,0.3) -- (0.3,0.3) -- (0.3,0);
\node at (-0.3,3.7) {$B$};
\node[below left] at (C) {$C$};
\node[below right] at (A) {$A$};
\pic[draw, angle radius=0.6cm, "$40^\circ$", angle eccentricity=1.5] {angle = C--B--A};
\end{tikzpicture}
\end{minipage}

\titreexo{Question 9}
Le diagramme en barres ci-dessous donne les notes des élèves d'une classe au dernier contrôle de mathématiques.

\begin{enumerate}[label=\textbf{\alph*.}]
\item Combien d'élèves ont participé à ce contrôle ?
\item Quelle est la note médiane ?
\end{enumerate}

\begin{center}
%Début Tikz
%\begin{tikzpicture}
%\begin{axis}[
%  ybar,
%  width=15cm,
%  height=6cm,
%  bar width=12pt,
%  xlabel={notes},
%  ylabel={effectif},
%  symbolic x coords={7,8,9,10,11,12,13,14,15,16,17,18},
%  xtick=data,
%  ymin=0, ymax=6,
%  ytick={0,1,2,3,4,5,6},
%  enlarge x limits=0.05,
%  axis lines=left,
%  nodes near coords,
%]
%\addplot coordinates {(7,3) (8,4) (10,4) (11,5) (12,5) (15,3) (17,2) (18,1)};
%\end{axis}
%\end{tikzpicture}
%Fin Tikz
% Début PSTricks
\psset{unit=0.9cm,arrowsize=2pt 3}
\begin{pspicture}(-1,-1)(12.5,6.5)
\psaxes[linewidth=1.25pt,Ox=6]{->}(0,0)(0,0)(12.5,6.5)
\psframe(0.8,0)(1.2,3)\psframe(1.8,0)(2.2,4)\psframe(3.8,0)(4.2,4)
\psframe(4.8,0)(5.2,5)\psframe(5.8,0)(6.2,5)\psframe(8.8,0)(9.2,3)
\psframe(10.8,0)(11.2,2)\psframe(11.8,0)(12.2,1)
\rput(6,-0.8){notes}\rput{90}(-0.8,3.25){effectifs}
\end{pspicture}
%Fin PSTricks
\end{center}

\begin{center}
\textit{Restitution de la copie du candidat à l'issue de la partie 1}
\end{center}

\section{Partie 2 -- Raisonnement et résolution de problèmes -- 14 points -- 1 h 40}

Dans cette partie, toutes les réponses doivent être justifiées, sauf si une indication contraire est donnée.

La clarté et la précision des raisonnements ainsi que la rédaction sont évaluées sur 2 points.

Pour chaque question, si le travail n'est pas terminé, laisser tout de même une trace de la recherche~; les essais et les démarches engagées, même non aboutis, seront pris en compte dans la notation.

\medskip

\textbf{Exercice 1 \hfill 2,5 points}

Lola souhaite acheter un smartphone. Elle étudie deux propositions.

\begin{center}
\begin{tabular}{p{7cm} p{7cm}}
\textbf{Offre A} & \textbf{Offre B} \\
Le client paie 175~\euro{} à l'achat, puis son abonnement est de 16~\euro{} par mois avec un engagement de 24 mois minimum. &
Le client ne paie rien à l'achat, puis l'abonnement est de 23~\euro{} par mois avec un engagement de 24 mois minimum.
\end{tabular}
\end{center}

\begin{enumerate}
\item Au bout de 24 mois, laquelle des deux offres est la plus intéressante ?
\item $x$ est un nombre positif qui représente le nombre de mois. On exprime le prix de ces deux tarifs en fonction de $x$, avec les fonctions suivantes :

\begin{itemize}
\item $f(x) = 175 + 16x$
\item $g(x) = 23x$
\end{itemize}

	\begin{enumerate}[label={\textbf{\alph*.}}]
		\item Associer chaque fonction à l'offre correspondante (A ou B).

\emph{Aucune justification n'est attendue.}
		\item Au bout de combien de mois paie-t-on le même prix avec ces deux offres ?
		\item Est-on encore dans la période d'engagement ?
	\end{enumerate}
\end{enumerate}

\textbf{Exercice 2 \hfill 3 points}

\emph{La figure ci-dessous n'est pas à l'échelle.}

\begin{minipage}{0.48\linewidth}
\begin{itemize}
\item O, A et B sont alignés.
\item O, D et C sont alignés.
\item OD $= 8,2$ cm
\item AD $= 1,8$ cm
\item BC $= 4,5$ cm
\end{itemize}
\end{minipage}\hfill
\begin{minipage}{0.48\linewidth}
% Début PSTricks
\psset{unit=1cm}
\begin{pspicture}(-1,-0.5)(3.6,5)
%\psgrid
\pspolygon(0,0)(0,5)(2.6,5)%ABCO
\psline(0,3)(1.6,3)%AD
\psframe(0,3)(0.3,2.7,3)
\psframe(0,5)(0.3,4.7)
\uput[d](0,0){O} \uput[l](0,3){A} \uput[l](0,5){B} \uput[r](2.6,5){C} \uput[r](1.6,3){D} 
\end{pspicture}
%Fin  PSTricks
\end{minipage}
%\begin{tikzpicture}[scale=1.05]
%% Points based on description: O at bottom, A,B aligned ; D,C aligned
%\coordinate (O) at (0,0);
%\coordinate (A) at (0.7,3.6);
%\coordinate (D) at (1.55,3.6);
%\coordinate (B) at (1.0,5.0);
%\coordinate (C) at (2.5,5.0);
%
%\draw (O) -- (B);
%\draw (O) -- (C);
%\draw (A) -- (D);
%\draw (B) -- (C);
%
%% right angle marks
%\draw ($(A)+(0.18,0)$) -- ($(A)+(0.18,0.18)$) -- ($(A)+(0,0.18)$);
%\draw ($(B)+(0.18,0)$) -- ($(B)+(0.18,0.18)$) -- ($(B)+(0,0.18)$);
%
%\node[left] at (O) {$O$};
%\node[left] at (A) {$A$};
%\node[right] at (D) {$D$};
%\node[above left] at (B) {$B$};
%\node[above right] at (C) {$C$};
%\end{tikzpicture}

\medskip

\begin{enumerate}
\item Montrer que la longueur du segment [OA] est égale à 8 cm.
\item Justifier que les droites (BC) et (AD) sont parallèles.
\item Calculer la longueur du segment [OB].
\item Une entreprise souhaite fabriquer des gobelets. Un gobelet (grisé sur le schéma ci-dessous) a la forme d'un tronc de cône (cône coupé par un plan parallèle à sa base).

\medskip

\begin{minipage}{0.48\linewidth}
On reprend les données précédentes :

\begin{itemize}
\item O, A, B sont alignés
\item O, D, C sont alignés
\item OD $= 8,2$ cm
\item AD $= 1,8$ cm
\item BC $= 4,5$ cm
\end{itemize}
\end{minipage}\hfill
\begin{minipage}{0.48\linewidth}
\emph{La figure ci-dessous n'est pas à l'échelle.}

\medskip

%Début Tikz
%\begin{center}
%\begin{tikzpicture}[scale=0.9]
%% big cone with frustum shaded
%\coordinate (O) at (0,0);
%\coordinate (Atop) at (-1.0,3.2); % small ellipse left
%\coordinate (Btop) at (-2.4,4.6); % large ellipse left
%
%% Outer cone lines
%\draw[dashed] (O) -- (-2.4,4.6);
%\draw[dashed] (O) -- (2.4,4.6);
%\draw[dashed] (O) -- (-1.0,3.2);
%\draw[dashed] (O) -- (1.0,3.2);
%
%% Large ellipse (top, cup rim)
%\draw (-2.4,4.6) ellipse (2.4 and 0.5);
%% small ellipse (base of cup)
%\draw (-1.0,3.2) ellipse (1.0 and 0.22);
%
%% Shaded frustum sides
%\filldraw[fill=gray!25, draw=black] (-1.0,3.2) -- (-2.4,4.1) arc (180:0:2.4 and 0.5) -- (1.0,3.2) arc(0:180:1.0 and 0.22) -- cycle;
%
%\node[right] at (2.4,4.6) {$C$};
%\node[left] at (-2.4,4.6) {$B$};
%\node[right] at (1.0,3.2) {$D$};
%\node[left] at (-1.0,3.2) {$A$};
%\node[below] at (O) {$O$};
%
%\draw[->, thick] (3.6,3.8) -- (1.6,3.9);
%\node[right] at (3.6,3.8) {Le gobelet};
%\end{tikzpicture}
%\end{center}
%Fin Tikz
%Début PSTricks
\begin{pspicture}(-2.6,0)(2.6,5)
%\psgrid
\psellipse(0,4.8)(2.3,0.3)
\psline[linestyle=dashed](-2.34,4.8)(2.34,4.8)
\psellipticarc(0,2.3)(1.15,0.2){180}{360}
\psellipticarc[linestyle=dashed](0,2.3)(1.15,0.2){0}{180}
\psline[linestyle=dashed](0,0)(0,4.8)%OB
%\psline[linestyle=dashed](-2.3,5)(2.3,5)
\psline[linestyle=dashed](-1.15,2.3)(1.15,2.3)
\psline(-2.34,4.8)(-1.13,2.3)
\psline(2.34,4.8)(1.13,2.3)
\psframe(0,4.8)(0.25,4.55)\psframe(0,2.3)(0.25,2.55)
\psline[linestyle=dashed](-1.15,2.3)(0,0)(1.15,2.3)
\uput[d](0,0){O} \uput[ul](0,2.3){A} \uput[ul](0,4.8){B} \uput[r](2.3,4.8){C} \uput[r](1.15,2.3){D} 
\rput(0,3.6){Le gobelet}
\end{pspicture}
%Fin PSTricks
\end{minipage}

\medskip

\fbox{\parbox{0.9\textwidth}{\textbf{Rappel :}\\ Volume d'un cône de révolution $V = \pi \times R^2 \times H \div 3$\\ où $R$ est le rayon de la base et $H$ est la hauteur du cône.}}

\begin{enumerate}[label=\textbf{\alph*.}]
\item Calculer le volume du grand cône de hauteur [OB] en cm$^3$, arrondi à l'unité.
\item Calculer le volume du gobelet, en cm$^3$, arrondi à l'unité.
\end{enumerate}
\end{enumerate}

\newpage

\textbf{Exercice 3 \hfill 4 points}

On considère le \textbf{programme A} suivant :

\begin{center}
\begin{tikzpicture}[
  box/.style={draw, rounded corners, align=center, minimum width=3.6cm, minimum height=0.9cm},
  >={Stealth[length=2.5mm]}
]
\node[box] (start) at (3,6) {Choisir un nombre};
\node[box] (carre) at (0,4.5) {Calculer son carré};
\node[box] (mult7) at (6,4.5) {Multiplier par 7};
\node[box] (mult3) at (0,3) {Multiplier par 3};
\node[box] (add) at (3,1.5) {Additionner les deux nombres};
\node[box] (sous) at (3,0) {Soustraire 6};

\draw[->] (start) -- (carre);
\draw[->] (start) -- (mult7);
\draw[->] (carre) -- (mult3);
\draw[->] (mult3) -- (add);
\draw[->] (mult7) -- (add);
\draw[->] (add) -- (sous);
\end{tikzpicture}
\end{center}

\begin{enumerate}
\item Appliquer le \textbf{programme A} au nombre 5.

\item On utilise un tableur pour trouver les résultats correspondants à quelques nombres comme l'indique le tableau ci-contre.

Parmi les quatre formules ci-dessous, recopier celle qui a été saisie dans la cellule B2, puis étirée vers le bas afin de calculer les résultats donnés par le programme A.

\emph{Aucune justification n'est attendue.}

\begin{center}
\begin{tabular}{|l|l|}\hline
\texttt{= 3*A2*2+7*A2-6} & \texttt{= 3*1*1+7*1-6} \\\hline
\texttt{= 3*A2*A2+7*A2-6} & \texttt{= 3*A2*2-7*A2+6} \\\hline
\end{tabular}
\end{center}

\begin{center}
\begin{tabular}{|c|c|c|}
\hline
 & \textbf{A} & \textbf{B} \\
\hline
\textbf{1} & Nombre de départ & Résultat du programme A \\
\hline
2 & $-3,5$ & 6,25 \\
3 & $-3$ & 0 \\
4 & $-2,5$ & $-4,75$ \\
5 & $-2$ & $-8$ \\
6 & $-1,5$ & $-9,75$ \\
7 & $-1$ & $-10$ \\
8 & $-0,5$ & $-8,75$ \\
9 & $0$ & $-6$ \\
10 & $0,5$ & $-1,75$ \\
11 & $1$ & $4$ \\
12 & $1,5$ & $11,25$ \\
13 & $2$ & $20$ \\
\hline
\end{tabular}
\end{center}

\item \`A l'aide du tableur, donner une valeur pour laquelle le \textbf{programme A} donne 0. \textit{Aucune justification n'est attendue.}

\item Si on note $x$ le nombre de départ, donner une expression littérale du programme A en fonction de $x$.
\end{enumerate}

\newpage

On considère maintenant le \textbf{programme B} suivant :

\begin{center}
\begin{tikzpicture}[
  box/.style={draw, rounded corners, align=center, minimum width=3.6cm, minimum height=0.9cm},
  >={Stealth[length=2.5mm]}
]
\node[box] (start) at (3,5) {Choisir un nombre};
\node[box] (mult3) at (0,3.5) {Multiplier par 3};
\node[box] (add3) at (6,3.5) {Additionner 3};
\node[box] (sous2) at (0,2) {Soustraire 2};
\node[box] (multdeux) at (3,0.5) {Multiplier les deux nombres};

\draw[->] (start) -- (mult3);
\draw[->] (start) -- (add3);
\draw[->] (mult3) -- (sous2);
\draw[->] (sous2) -- (multdeux);
\draw[->] (add3) -- (multdeux);
\end{tikzpicture}
\end{center}

\begin{enumerate}[resume]
\item Appliquer le \textbf{programme B} au nombre 5.

\item Si on note $x$ le nombre de départ, donner une expression littérale du \textbf{programme B} en fonction de $x$.

\item Mathis affirme que, quel que soit le nombre qu'il choisit, il trouvera le même résultat avec le \textbf{programme A} et le \textbf{programme B}. A-t-il raison ? Justifier.

\item Résoudre l'équation $(3x-2)(x+3) = 0$.

En déduire les valeurs de $x$ pour lesquelles les programmes A et B donnent 0.
\end{enumerate}

\bigskip

\textbf{Exercice 4 \hfill 2,5 points}

\begin{minipage}{0.48\linewidth}
ABCD est un carré.

ABF est un triangle équilatéral.

AF $= 4$ cm.

Les points F{}, A et E sont alignés.
\end{minipage}\hfill
\begin{minipage}{0.48\linewidth}
\psset{unit=1cm}
\begin{pspicture}(-3.5,-1.4)(3.5,4.8)
%\psgrid
\psline(-3.2,0)(3.2,0)
\pspolygon(0;0)(3.2;120)(3.2;180)
%\rput{30}(0,0){\psframe(0,0)(3.2,3.2)}
\def\barb{\psline(-0.05,-0.1)(-0.05,0.1)\psline(0.05,-0.1)(0.05,0.1)}
\rput{30}(0,0){\psline(0;0)(3.2;0)(4.525;45)(3.2;90)}
\psarc*(0,0){0.6}{0}{30}
\uput[d](1.6;180){4 cm}\uput[d](0,0){A} \uput[d](3.2,0){E}\uput[ul](-1.6,2.772){B}\uput[r](3.2;30){D}\uput[ur](1.4,4){C}\uput[d](-3.2,0){F}
\rput(-1.6,0){\barb}\rput(1.6,0){\barb}
\rput{30}(1.6;30){\barb}\rput{135}(2,3){\barb}\rput{120}(-0.8,1.4){\barb}
\rput{600}(-2.4,1.4){\barb}\rput{30}(-0.2,3.6){\barb}
\psdots[dotstyle=+,dotangle=45,dotscale=1.5](3.2,0)
\rput(0,-1){\textit{Cette figure n'est pas en vraie grandeur.}}
\end{pspicture}

\end{minipage}
%Début Tikz
%\begin{center}
%\begin{tikzpicture}[scale=0.85]
%\coordinate (F) at (0,0);
%\coordinate (A) at (4,0);
%\coordinate (E) at (7,0);
%\coordinate (B) at (2,3.46);
%\coordinate (C) at (5.46,5.46);
%\coordinate (D) at (7.46,2);
%
%\draw (F) -- (A) -- (B) -- cycle;
%\draw (A) -- (B) -- (C) -- (D) -- cycle;
%\draw[->] (A) -- (E);
%\draw (F) -- (A);
%
%% tick marks for equal sides FA, AB, BF
%\draw (1.9,0.15) -- (2.1,-0.15);
%\draw (3.0,1.85) -- (3.2,1.6);
%\draw (0.95,1.85) -- (0.75,1.6);
%
%% small square tick marks for ABCD sides
%\draw (4.7,0.05)--(4.9,-0.25);
%\draw (6.46,4.5)--(6.66,4.2);
%\draw (3.7,4.5)--(3.9,4.8);
%
%\node[left] at (F) {$F$};
%\node[below] at (A) {$A$};
%\node[right] at (E) {$E$};
%\node[above left] at (B) {$B$};
%\node[above] at (C) {$C$};
%\node[right] at (D) {$D$};
%\node[below] at (3,-0.3) {4 cm};
%
%\pic[draw, green!50!black, angle radius=0.5cm, angle eccentricity=1.6] {angle = E--A--D};
%\end{tikzpicture}
%\end{center}
%Fin Tikz

\begin{center}

\end{center}

\begin{enumerate}
\item Justifier que la mesure de l'angle $\widehat{\text{EAD}}$ est égale à $30^\circ$.

Dans la suite de l'exercice on utilisera l'échelle suivante : \textbf{10 pas} dans le programme représentent \textbf{1 cm} dans la réalité.

\item Le \textbf{bloc triangle} ci-dessous permet de tracer un triangle équilatéral et le \textbf{bloc carré} permet de construire un carré :

\begin{center}
\begin{tabular}{|p{7cm}|p{7cm}|}
\hline
\textbf{Le bloc triangle} & \textbf{Le bloc carré} \\
\hline
%\begin{tikzpicture}[scale=0.8]
%\node[draw, fill=red!60, rounded corners, minimum width=3cm] at (0,3) {définir Triangle};
%\node[draw, fill=orange!70, rounded corners, minimum width=3cm] (r1) at (0,2) {répéter 3 fois};
%\node[draw, fill=blue!40, rounded corners, minimum width=3cm] at (0.3,1.2) {avancer de \textbf{J} pas};
%\node[draw, fill=blue!40, rounded corners, minimum width=3cm] at (0.3,0.4) {tourner de \textbf{K} degrés};
%\end{tikzpicture}
\begin{scratch}
\initmoreblocks{définir \namemoreblocks{Triangle}}
\blockrepeat{répéter \ovalnum{3} fois}
{
\blockmove{avancer de \ovalnum{\textbf{J}}}
\blockmove{tourner \turnleft{} de \ovalnum{\textbf{K}} degrés}
}
\end{scratch}
&
\begin{scratch}
\initmoreblocks{définir \namemoreblocks{Carré}}
\blockrepeat{répéter \ovalnum{4} fois}
{
\blockmove{avancer de \ovalnum{\textbf{M}}}
\blockmove{tourner \turnleft{} de \ovalnum{N} degrés}
}
\end{scratch}
%\begin{tikzpicture}[scale=0.8]
%\node[draw, fill=red!60, rounded corners, minimum width=3cm] at (0,3) {définir Carré};
%\node[draw, fill=orange!70, rounded corners, minimum width=3cm] (r1) at (0,2) {répéter 4 fois};
%\node[draw, fill=blue!40, rounded corners, minimum width=3cm] at (0.3,1.2) {avancer de \textbf{M} pas};
%\node[draw, fill=blue!40, rounded corners, minimum width=3cm] at (0.3,0.4) {tourner de \textbf{N} degrés};
%\end{tikzpicture}\\
\\ \hline
\end{tabular}
\end{center}

Donner les valeurs de \textbf{J}, \textbf{K}, \textbf{M} et \textbf{N} pour que les blocs triangle et carré permettent de construire un triangle équilatéral de 4 cm de côté et un carré de 4 cm de côté.

\emph{Aucune justification n'est attendue.}

\item Le \textbf{programme principal} utilise le bloc \textbf{Triangle} et le bloc \textbf{Carré}.

L'instruction \og s'orienter à $90\degres$ \fg{} signifie que l'on s'oriente vers la droite.

Écrire sur la copie le numéro de la figure obtenue grâce à ce programme.

\emph{Aucune justification n'est attendue.}

\noindent
%\begin{minipage}[t]{0.42\textwidth}
%\textbf{Figure 1 :}\\[0.3em]
%\begin{tikzpicture}[scale=0.4]
%\foreach \x in {0,...,7} {
%  \draw (\x*1.4,0) -- (\x*1.4+0.5,1) -- (\x*1.4+1,0);
%  \draw (\x*1.4+1,0) rectangle (\x*1.4+1.4,0.5);
%}
%\end{tikzpicture}
%\end{minipage}%
%\hfill
%\begin{minipage}[t]{0.55\textwidth}
%\textbf{Programme principal}\\[0.3em]
%\begin{tikzpicture}[
%  box/.style={draw, rounded corners, align=center, minimum width=4.4cm, minimum height=0.7cm, font=\small},
%  >={Stealth[length=2mm]}
%]
%\node[box, fill=yellow!60] (q1) at (0,7.4) {quand le drapeau est cliqué};
%\node[box, fill=blue!30] (q2) at (0,6.6) {aller à x : 0 \ y : -100};
%\node[box, fill=blue!30] (q3) at (0,5.8) {s'orienter à 90};
%\node[box, fill=green!40] (q4) at (0,5.0) {effacer tout};
%\node[box, fill=green!40] (q5) at (0,4.2) {stylo en position d'écriture};
%\node[box, fill=orange!60] (q6) at (0,3.4) {répéter 6 fois};
%\node[box, fill=red!50] (q7) at (0.3,2.5) {Triangle};
%\node[box, fill=blue!30] (q8) at (0.3,1.75) {avancer de 50 pas};
%\node[box, fill=blue!30] (q9) at (0.3,1.0) {tourner de 30 degrés};
%\node[box, fill=red!50] (q10) at (0.3,0.25) {Carré};
%\node[box, fill=blue!30] (q11) at (0.3,-0.5) {avancer de 50 pas};
%\node[box, fill=blue!30] (q12) at (0.3,-1.25) {tourner de 30 degrés};
%\draw[->] (q1) -- (q2);
%\draw[->] (q2) -- (q3);
%\draw[->] (q3) -- (q4);
%\draw[->] (q4) -- (q5);
%\draw[->] (q5) -- (q6);
%\draw[->] (q6) -- (q7);
%\draw[->] (q7) -- (q8);
%\draw[->] (q8) -- (q9);
%\draw[->] (q9) -- (q10);
%\draw[->] (q10) -- (q11);
%\draw[->] (q11) -- (q12);
%\end{tikzpicture}
%\end{minipage}

%\vspace{1em}
%
%\noindent
%\begin{minipage}[t]{0.45\textwidth}
%\textbf{Figure 2 :}\\[0.3em]
%\begin{tikzpicture}[scale=0.5]
%\draw (0,0) circle (2);
%\foreach \i in {0,...,7} {
%\draw ({2*cos(45*\i)},{2*sin(45*\i)}) -- ({1.4*cos(45*\i+22.5)},{1.4*sin(45*\i+22.5)}) -- ({2*cos(45*\i+45)},{2*sin(45*\i+45)});
%}
%\end{tikzpicture}
%\end{minipage}%
%\hfill
%\begin{minipage}[t]{0.45\textwidth}
%\textbf{Figure 3 :}\\[0.3em]
%\begin{tikzpicture}[scale=0.5]
%\draw (0,0) circle (2);
%\foreach \i in {0,...,7} {
%  \draw ({2*cos(45*\i)},{2*sin(45*\i)}) rectangle ++(0.6,0.6);
%}
%\end{tikzpicture}
%\end{minipage}

%\vspace{1em}
%
%\noindent
%\textbf{Figure 4 :}\\[0.3em]
%\begin{tikzpicture}[scale=0.4]
%\foreach \x in {0,...,5} {
%\draw (\x*1.6,0) -- (\x*1.6+0.5,1) -- (\x*1.6+1,0);
%\draw (\x*1.6+1,0) rectangle (\x*1.6+1.4,0.5);
%}
%\end{tikzpicture}


\end{enumerate}
%PSTricks
\begin{tabularx}{\linewidth}{|m{8.9cm}|X|}\hline
\textbf{Figure 1}&\textbf{Programme principal}\\
\begin{tabular}{|m{8.5cm}|}
\hline
\psset{unit=0.8cm}
\quad\begin{pspicture}(0,-0.5)(8.5,1.5)
\def\motif{\pspolygon[linecolor=blue](0;0)(0.75;120)(0.75;180)\psframe[linecolor=blue](0,0)(0.75,0.75)}
\multido{\n=0.75+1.50}{6}{\rput(\n,0){\motif}}
\end{pspicture}\\ \hline
\end{tabular}

\begin{tabular}{|m{3.75cm}|m{4.25cm}|}\hline
\textbf{Figure 2}&\textbf{Figure 3}\\
\begin{center}
{\psset{unit=0.7cm}
\begin{pspicture}(-2,-2)(2,2)
%\psgrid
\def\dim{1}%%% modifiable 
\def\motif{\psline[linecolor=blue,linewidth=1.2pt](0,0)(\dim,0)(\dim,-\dim)(0,-\dim)(0,0)(\dim;60)(\dim,0)}
\multido{\n=30+60,\na=0+60}{6}{\rput{\na}(\dim;\n){\motif}}
\end{pspicture}}\end{center}
&
\begin{center}
{\psset{unit=0.9cm}
\begin{pspicture}(-2,-2)(2,2)
%\psgrid
\def\motif{\psline[linecolor=blue,linewidth=1.2pt](0,0)(0,-1)(-0.8,-1)(-0.8,-0.2)(0,-0.2)(0,0)}
\def\motiff{\psline[linecolor=blue,linewidth=1.2pt](0,0)(0.8;120)(0.8;180)(0,0)(1;120)}
\multido{\n=15+60.0,\na=0+60}{6}{\rput{\na}(1.9;\n){\motif}}
\multido{\n=15+60.0,\na=0+60}{6}{\rput{\na}(1.9;\n){\motiff}}
\end{pspicture}}
\end{center}\\ \hline
\end{tabular}

\textbf{Figure 4}

{\psset{unit=0.8}
\quad\begin{pspicture}(0,-0.5)(9,1.5)
\def\motif{\pspolygon[linecolor=blue](0;0)(0.75;120)(0.75;180)\psframe[linecolor=blue](0,0)(0.75,0.75)\psline(0.75,0)(1,0)}
\multido{\n=1.00+1.75}{6}{\rput(\n,0){\motif}}
\end{pspicture}}
&
\vspace{-4cm}
\begin{scratch}[scale=0.85]
\blockinit{quand \greenflag est cliqué}
\blockmove{aller à x : \ovalnum{0} y : \ovalnum{-100}}
\blockmove{s'orienter à ~\ovalnum{90}}
\blockpen{effacer tout}
\blockpen{stylo en position d'écriture}
\blockrepeat{répéter \ovalnum{6} fois}
{
\blockmoreblocks{triangle}
\blockmove{avancer de \ovalnum{50} pas}
\blockmove{tourner \turnleft{} de \ovalnum{30} degrés}
\blockmoreblocks{Carré}
\blockmove{avancer de \ovalnum{50} pas}
\blockmove{tourner \turnleft{} de \ovalnum{30} degrés}
}
\end{scratch}\\ 
\hline
\end{tabularx}
%Fin PSTricks
\end{document}